\chapter{Discussion}\label{chap:key:7}

The study aimed to attest whether the word order in the Turkish and Kurmanji varieties is undergoing change, with a specific focus on the post-predicate domain. The two languages are different in the context of heritage language research: while Turkish is one of the most vital heritage languages and has been extensively studied at least in the context of Germany, Kurmanji remains an underexplored heritage language, particularly in the German context.

This study analyzed post-predicate structures in Turkish across four groups: two Turkish heritage speaker groups in Germany—Turkish-German bilinguals and Turkish-Kurmanji-German trilinguals—and two majority Turkish speaker groups in Turkey—Turkish-Kurmanji bilinguals and Turkish monolinguals. The rich dataset for Turkish provided a valuable opportunity to investigate the impact of various factors on the use of the post-predicate position, several of which have not been thoroughly examined in the existing literature.

First, this study shifts the focus from merely comparing bilinguals to monolinguals, instead highlighting the role of societal language status—heritage versus majority. In this context, majority status includes both bilingual and monolingual speakers, while heritage status encompasses both bilinguals and trilinguals. Second, different contact situations are considered, taking into account the various contact languages involved: Kurmanji in Turkey, German in Germany, and both languages among trilingual speakers. Additionally, the intensity of contact varies across the groups. For instance, the contact between Turkish and Kurmanji began before the Ottoman Empire and intensified following the establishment of the Turkish Republic (\citealt{Yagmur2001}). In contrast, contact between Kurmanji and German emerged in the second half of the 20th century with the arrival of guest workers in Germany (\citealt{GhaderiAlmstadt2023}). Finally, as emphasized throughout this study, it is crucial to consider register when investigating heritage languages. Many phenomena examined in heritage language studies are susceptible to register variation, with certain patterns predominantly observed in formal registers. Therefore, it is essential to compare heritage speakers with appropriate data from majority speakers. In other words, comparing heritage language use in informal settings with monolingual productions in formal settings would lead to an unfair comparison that disadvantages bilingual speakers.

In the context of investigating word order in Turkish, it is particularly important to account for register variation. The use of the post-predicate position is not noncanonical per se, instead, it is typical of informal language use, while its occurrence is restricted in formal registers. Taking these arguments into account, the data for this study were collected using a method that captures the broader repertoires of speaker groups across both formal and informal settings, as well as spoken and written modes \citep{Wiese2020}. This approach was applied to all speaker groups, yielding data used across four different communicative situations among both heritage and majority speakers.

Unlike Turkish, the impact of various factors on Kurmanji cannot be thoroughly investigated. First, assessing the impact of the societal status of Kurmanji poses challenges. This study focuses on two speaker groups: Kurmanji-Turkish bilinguals in Turkey and Kurmanji-Turkish-German trilinguals in Germany. In both cases, Kurmanji is classified as a heritage language. Second, there is a scarcity of monolingual Kurmanji speakers, as most individuals are typically bilingual or multilingual, often demonstrating greater proficiency in the majority language of the country in which they reside, such as Turkish, Persian, Arabic, or other contact languages (\citealt{BrizicnodateEtAl}). Furthermore, although Kurmanji is studied in two countries with different majority languages, Turkish serves as the contact language for both groups. However, Turkish holds different societal statuses in each context: it is a majority language in Turkey and a heritage language in Germany. Additionally, the intensity of contact varies between the two language constellations, with the contact between Kurmanji and Turkish in Turkey having a much longer history than that between Kurmanji and German in Germany. Finally, this study also examines the role of register in word order variation in Kurmanji. Unlike Turkish, the use of the post-predicate position in Kurmanji is primarily influenced by predicate semantics and is not anticipated to be affected by register. However, several factors justify investigating the impact of register. First, there has been no research to date on the influence of formality and mode on word order in Kurmanji. Additionally, studies have shown how formality and mode affect word order in Persian, a language that shares syntactic features with other Western Iranian languages, including Kurmanji (\citealt{Frommer1981}, \citealt{Haig_post-predicate_EtAl2024-1}). Furthermore, in Turkish, which serves as the contact language for both groups of Kurmanji speakers in this study, the post-predicate position is related to formality and mode. Finally, as a heritage language, Kurmanji is typically not used in formal situations. Nevertheless, since speakers are likely aware of how formality and mode influence language use in the majority language, it is interesting to explore whether they develop new strategies in their heritage language to adhere to the linguistic conventions of specific communicative situations.

The main results and the underlying causes for the observed findings, along with their relation to previous studies on Turkish and Kurmanji, have been summarized at the end of Chapters 5 and 6, respectively. Given the numerous factors considered and their interplay, the main findings will be summarized in three sections: \sectref{sec:key:7.1} discusses evidence of CLI between the languages. \sectref{sec:key:7.2} explores the implications of the results for understanding internal linguistic dynamics, while \sectref{sec:key:7.3} highlights findings related to register. Finally, \sectref{sec:key:7.4} outlines the study’s limitations.

\section{Evidence of CLI on word order in Turkish and Kurmanji}\label{sec:key:7.1}

The languages under investigation here, Turkish, Kurmanji, and German, are SOV languages, but they employ the post-predicate position to a different extent. Moreover, the employment of this position is determined by different factors in the three languages: information structure in Turkish (\citealt{Erguvanli1984, Erku1983, Issever2003}), predicate semantics in Kurmanji (\citealt{Gundogdu2019, Haig2018, Haig2022}), and syntactic structure in German (\citealt{Frey2015}, \citealt{Muller2023}).

At the same time, there are also overlaps in the factors that influence the employment of the post-predicate position. Thus, for example, \citet{Asadpour2021} suggests that in addition to predicate semantics, information structure also impacts the use of the post-predicate position in Kurdish. If information structure similarly has an effect on the word order of Kurmanji spoken in Turkey, it would highlight an overlapping factor with Turkish -- information structure. Furthermore, common factors for the employment of the post-predicate position in Turkish and German are mode and formality: the post-predicate position occurs more frequently in the spoken mode compared to the written mode and in the informal language compared to the formal one (\citealt{Erguvanli1984, Schroeder2002}, \citealt{Tsehaye2023}).

Thus, the interplay of similarities and differences in these factors renders the investigated topic particularly interesting for examination in terms of CLI. This study aimed to determine whether the use of the post-predicate position varies across different language constellations in terms of frequency distribution and whether innovative factors influence its use in language contact situations. CLI refers to instances where there is evidence of one language affecting another. The societal status of the languages plays an important role in defining the direction of influence: usually a more prestigious and dominant language influences the less prestigious one. Thus, studies that investigate contact-induced language change often concentrate on minority languages. This study takes a broader approach. First, it examines Turkish in contact settings where it functions both as a heritage and a majority language. Second, it investigates potential word order changes in Kurmanji, a language in contact with Turkish. This dual focus provides deeper insights into the sources of structural changes in both languages. This section first presents evidence for CLI on heritage Turkish and heritage Kurmanji before turning to the findings on majority Turkish.

Regarding the Turkish heritage groups in Germany, the frequency analysis of post-predicate structures presented in Chapter \ref{chap:key:5} yielded mixed results. While Turkish-German bilingual speakers exhibit a frequency of post-predicate structures comparable to that of monolinguals, Turkish-Kurmanji-German trilingual speakers use them less frequently. However, a similar pattern is observed among Kurmanji-Turkish bilinguals who do not have German as a contact language, and notably, this difference emerges due to register inflexibility (see \sectref{sec:key:7.3} for further discussion). Additionally, the qualitative analysis of constituents in the post-predicate position revealed no differences between the groups. Therefore, the findings on post-predicate structures in heritage Turkish—whether among Turkish-German bilinguals or Kurmanji-Turkish-German trilinguals—do not point to the CLI from German.

With respect to Kurmanji, another heritage language under investigation, detecting and proving CLI is more challenging compared to Turkish, as monolingual Kurmanji speakers are virtually nonexistent, with most being at least bilingual. Furthermore, in both groups examined, Kurmanji is in contact with Turkish, making it difficult to establish a baseline for comparison. Therefore, as a reference point, this study draws on findings from previous research on word order in Kurmanji (\citealt{Haig2014}, \citealt{Haig2018, Gundogdu2019}), as well as studies on other Kurdish languages \citep{Asadpour2021} and Persian (\citealt{Frommer1981}, \citealt{Haig_post-predicate_EtAl2024-1}).

\largerpage
First, regarding the frequency of post-predicate structures across the groups, the analysis revealed that Kurmanji speakers in Germany use these structures more frequently than those in Turkey. However, this difference is primarily due to the more frequent placement of goals of motion in the post-predicate position among the former group. Since the post-predicate position is the canonical position for goals in Kurmanji, this distribution does not provide direct evidence for CLI from either German in Germany or Turkish in Turkey. Instead, the observed difference may simply reflect a more frequent expression of motion by the speakers in Germany.

In contrast, the qualitative analysis of goal arguments in Kurmanji indicated a structural change at the level of constructions, suggesting that word order changes are observed only in particular constructions. While the canonical position for goal arguments flagged with case or a locational noun is immediately after the predicate, there are examples in the data where such goals are placed in the pre-predicate position. This noncanonical placement occurs predominantly in certain constructions.

In Chapter \ref{chap:key:6}, it was discussed that a possible reason for the word order change at the level of constructions in Kurmanji is the influence of Turkish. One example is the construction \textit{erdê ketin} (‘to fall on the ground’), where the case-flagged argument \textit{erdê} (‘on the ground’) is placed before the predicate. This placement may be a result of the specific Turkish pattern \textit{yere düşmek}. However, there is no similar evidence for other constructions, such as \textit{hemberî rê çun/revîn} (‘to go/run to the opposite side of the road’). Thus, this suggests that the preverbal position of goal arguments marked with a locational noun in Kurmanji may be more accurately attributed to the copying of the schematic [N-DAT V] pattern from Turkish.

Similarly, the Kurmanji data from Germany revealed instances of goals marked with case or locational nouns placed in the pre-predicate position. However, unlike Turkish, no evidence suggests that these noncanonical examples in the German data resulted from copying specific or schematic constructions from German. In German, the word order follows an SOV pattern, but certain constituents, such as prepositional phrases, determiner phrases, adverbials, and adjectives, can be extraposed. Furthermore, in main declarative clauses with a finite main verb, the word order pattern is V2. Thus, while the SOV pattern in German may prompt the noncanonical placement of goals in Kurmanji, the V2 pattern might simultaneously induce the use of canonical post-predicate placement of goals in Kurmanji in Germany.

\largerpage
The baseline for Kurmanji in Germany is Kurmanji in Turkey, indicating that heritage speakers in Germany acquire the language from their bilingual parents. Findings from the data in Turkey demonstrate that Kurmanji is undergoing changes at the level of constructions, particularly with word order, showing a tendency to place goal arguments in the pre-predicate position in certain constructions. Therefore, the noncanonical placement of goal arguments observed in the German data likely does not result from contact with German, rather, these constructions were already present in the language of the parents and were subsequently acquired by heritage speakers in Germany.\footnote{{Although, in order to prove that the noncanonical position of goals is already present in the input received by the heritage speakers of Kurmanji in Germany, it is necessary to collect data from the parents of the speakers.}}

At the same time, the absence of direct influence from German on Kurmanji syntax may be attributed to the focus on noncanonical constructions. The V2 structure in German could also induce the canonical [V N-DAT] pattern in Kurmanji. It is important to note that the frequency analysis revealed a higher number of goal arguments in the post-predicate position in the data from Germany compared to that from Turkey. This observation may indeed reflect the influence of the majority language, German, on the stability of Kurmanji syntax.

Furthermore, the systematic observation of noncanonical placement of goal arguments in the same constructions in the data of several speakers from Turkey indicates that this pattern is not merely an error. Instead, these innovative constructions appear to be in the process of propagation. Both constructions—those with arguments in the pre-predicate and post-predicate positions—are firmly established in individual speakers’ minds, and it remains unclear whether the innovative construction will ultimately supplant the canonical one.

In contrast, the number of noncanonical constructions in the data from Germany is limited, with only three examples in total, two of which originate from the same speaker. Consequently, there is no evidence of propagation for these innovative constructions in Germany. However, it is important to reiterate that the dataset from Germany is small. Therefore, to determine whether the innovative pattern is used by a wider range of speakers and whether it extends to additional constructions, further data collection from Kurmanji heritage speakers in Germany is necessary.

The results outlined above pertain to innovative developments in the post-predicate domain in heritage Turkish and heritage Kurmanji. In addition to examining these heritage languages, the study also explored Turkish as the majority language of Kurmanji-Turkish bilinguals and Turkish monolinguals in Turkey. The analysis revealed that Kurmanji-Turkish bilinguals use fewer post-predicate structures compared to Turkish monolinguals. Notably, a similar pattern in frequency distribution was observed among the Turkish data from the Kurmanji-Turkish-German trilingual group in Germany.

This difference in frequencies could suggest a CLI from Kurmanji, which is predominantly a verb-final language. However, a detailed analysis of post-predicate structures indicates that the frequency differences between the groups with and without Kurmanji as a contact language do not stem from such influence. Firstly, the qualitative analysis of post-predicate structures in both groups showed no evidence of CLI at the construction level. For example, verb arguments—such as goals of motion, addressees of verbs of speech, and recipients of verbs of transfer—that are systematically placed in the post-predicate position in Kurmanji do not occur more frequently in that position in Turkish among groups where Kurmanji serves as the contact language compared to those without it. Secondly, the analysis demonstrated that differences in the frequency distribution of post-predicate structures between the groups with and without Kurmanji as a contact language primarily result from register inflexibility (see \sectref{sec:key:7.3} for further discussion). If CLI from Kurmanji on Turkish were present, one would expect to see fewer post-predicate structures in groups where Kurmanji is the contact language across all communicative situations. However, since this is not the case, the observed differences in frequency cannot be attributed to the CLI.

In summary, the analysis found no evidence of CLI between the recently contacted languages, specifically heritage Turkish and Kurmanji in relation to German. Regarding the longer-standing contact between Turkish and Kurmanji, the societal status of the languages significantly influences the direction of that influence. The analysis did not reveal evidence of CLI in majority Turkish. However, it indicated that Kurmanji, as a heritage language, is experiencing CLI from Turkish, evidenced by changes in word order at the construction level.

According to the usage-based approach, change in syntax is the result of the change of many constructions, which are not spontaneous, but systematic changes. The data indicate that word order changes at the construction level are in the process of propagation in Turkey, as noncanonical constructions are used by different speakers. However, these novel constructions do not yet surpass canonical constructions in frequency, leaving it uncertain whether the noncanonical patterns will eventually dominate. In contrast, the data from Germany show only a few noncanonical constructions, suggesting no evidence of propagation for these innovative structures in that context. Nonetheless, any claims regarding Kurmanji in Germany remain tentative, as the dataset is too small to draw definitive conclusions.

\section{Internal dynamics: Is monolingual Turkish undergoing change?}\label{sec:key:7.2}

Apart from the contact situations, the investigation also focused on the monolingual variety, although only of Turkish. As previously explained, Kurmanji speakers are typically bilingual or multilingual, therefore, data from monolingual Kurmanji speakers are unavailable.

The analysis revealed several internal developments occurring in Turkish. First, the examination of the information status of referents in the post-predicate position indicated that monolingual speakers also place new referents in this position. This finding contradicts the assertions made in descriptive grammars of Turkish, which state that new information is not permitted in the post-predicate position (\citealt{Erguvanli1984, GokselÖzsoy2000, Issever2003}).

However, it was observed that most instances of new referents in the post-predicate position included secondary information about a nominal referent previously introduced in the pre-predicate position. Although research on this development in Turkish is limited, several studies have noted that new information can be placed in the post-predicate position if it is deemed less important compared to the new information introduced in the pre-predicate position (\citealt{DogruozBackus2007}, \citealt{AkanHartmann2019, Issever2019}). Additionally, the current study indicates that new elements in the post-predicate position occur in monolinguals, but primarily in informal contexts. This finding may suggest the initial stages of change within the language of monolingual speakers, however, it remains uncertain whether this will lead to a loosening of the information structural constraints in the future.

The second development concerning monolingual Turkish involves the use of post-predicate structures in formal situations. The study revealed that the post-predicate position is characteristic not only of informal settings, as stated by \citet{Erguvanli1984}, but is also frequently used in formal situations, particularly in spoken mode. It is important to note that this finding cannot be generalized across all formal spoken communicative situations in Turkish. The formal spoken production analyzed in this study comes from an oral witness report made to a police answering machine, where participants had to respond spontaneously without preparation. In contrast, more planned and prepared formal spoken situations, such as conference talks, news reports, or lectures, may contain significantly fewer post-predicate structures than an oral witness report. Future research is needed to gain a deeper understanding of the use of the post-predicate position across various formal registers.

\section{The role of register in heritage language research}\label{sec:key:7.3}

\largerpage
When comparing two or more speaker groups, selecting the appropriate baseline is crucial. This consideration is particularly important in the context of heritage languages, which are typically acquired at home in informal settings. Heritage speakers often have limited exposure to formal registers of their heritage language, as they frequently do not receive education in that language. Consequently, comparing spoken data collected from informal situations in a heritage language to the norms of a standard language may present a distorted view of developments in the heritage language \citep{WieseEtAl2022}. For instance, differences between the two may be mistakenly attributed to incomplete acquisition, attrition, or contact-induced language change. In reality, the observed differences between heritage and majority varieties might stem from register inflexibility in heritage speakers. The term \textit{register flexibility} is usually discussed in the context of foreign language learning, where register flexibility is defined as “the ability to flexibly use a variety of linguistic resources – at the lexical, syntactic and discourse levels, with the awareness of which are the most appropriate for the communicative contexts at hand” (\citealt{QinUccelli2020}).

Taking these considerations into account, the data for this study were collected using a method designed to capture the broader repertoires of speaker groups across both formal and informal settings in spoken and written modes \citep{Wiese2020}. As previously mentioned, \citet{SchroederEtAl2024} employed the same method to investigate the post-predicate position among Turkish heritage speakers in Germany and the US, finding evidence of register inflexibility in the latter group. Specifically, US speakers used the post-predicate position frequently in both informal and formal situations, whereas monolinguals typically employed this structure only in informal situations and spoken modes. However, no lack of register flexibility was identified among heritage speakers in Germany.

Consistent with the findings of \citet{SchroederEtAl2024}, the analysis of the current data did not reveal register inflexibility among Turkish heritage speakers in Germany in terms of the frequency distribution of post-predicate structures. However, this conclusion applies only to bilingual Turkish heritage speakers in Germany, who use significantly higher number of post-predicate structures in the informal situations compared to the formal ones, as well as in spoken mode versus written mode.

In terms of the frequency distribution of post-predicate structures, both overall and across various communicative situations, the data reveal differences between the groups with and without Kurmanji as the contact language. While Turkish monolinguals and Turkish heritage speakers in Germany show clear distinctions across communicative situations, Turkish bilingual speakers in Turkey show no significant differences regarding formality, and Kurmanji-Turkish heritage speakers in Germany exhibit no significant differences in either formality or mode. This pattern observed among Kurmanji speakers in both Germany and Turkey may suggest register inflexibility, where the distinctions between formal and informal situations, as well as between spoken and written language, become less pronounced.

First, the difference between the two heritage speaker groups in Germany is interesting: bilingual heritage speakers distinguish between different registers, while trilingual speakers do not. This variation may be attributed to their differing levels of exposure to and use of heritage Turkish. Bilingual speakers are exposed to Turkish at home and in informal public settings, particularly in Berlin, whereas trilingual speakers are also exposed to Kurmanji in addition to Turkish. However, establishing a correlation between the amount of input and exposure and the employment of post-predicate constituents in the different communicative situations falls outside the scope of this study.

Another contributing factor could be the different input received by the two groups: trilingual speakers acquire their heritage languages from their Kurmanji-Turkish bilingual parents. As the analysis has shown, the Turkish of Kurmanji-Turkish bilinguals is predominantly verb-final.  Furthermore, it is important to consider that the trilingual speaker group consists of only nine individuals, thus, increasing the sample size could potentially impact the observed frequency distribution.

At the same time, the difference between the two majority speaker groups merits attention. That is, monolingual majority speakers differentiate between formal and informal, as well as spoken and written communicative situations regarding their use of the post-predicate position. In contrast, Turkish bilingual speakers in Turkey do not exhibit significant differences based on formality. This frequency distribution suggests register inflexibility among bilingual majority speakers. However, this register inflexibility differs from what \citet{SchroederEtAl2024} observed in heritage speakers in the US, where Kurmanji-Turkish bilinguals use a greater number of structures typical of formal language in informal situations.

There are two potential explanations for the differences between the groups. First, Kurmanji-Turkish bilinguals in Turkey are largely exposed to Turkish in formal environments, while they primarily use Kurmanji with family members and some friends. As a result, the entrenchment of preverbal constructions is high among these speakers, influenced not only by the frequent SOV pattern in Kurmanji but also by their exposure to predominantly preverbal structures in formal Turkish situations.

Another reason may stem from the canonical behavior of bilingual speakers when using their majority language. Although research on the majority language use of heritage speakers is limited, evidence suggests that bilinguals tend to be more explicit and favor formal language and canonical structures when speaking their majority language (\citealt{SevincDewaele2018}, \citealt{Bunk2024, bunk_anxious_2025}). Studies indicate that monolingual habits and standard language ideologies can induce language anxiety in bilingual individuals, which may in turn lead to specific linguistic behaviors, such as the use of formal language in informal situations. As \citet{Bartholoma2021} indicates, language ideology in Turkey is often described as a `monoglot' ideology (for the term see \citealt{Silverstein1996}), which denies the existence of multilingualism within its borders and further restricts the use of non-state languages in both public and private spheres. This ideology embraces the notion of `one country, one language' and seeks to protect the state language from the influence of foreign and minority languages.

While there are no systematic studies on the societal perception of Turkish as spoken by bilingual speakers in Turkey, there is a recognized bias against Kurmanji-Turkish bilinguals, with the belief that their language does not conform to the norms of Standard Turkish. Additionally, several studies indicate that bilingual speakers in Turkey experience language anxiety (\citealt{Sahin-firat2010, Tulgar2010}). As a result, the use of formal markers and canonical language in the informal situations by bilinguals in Turkey might be the result of their wish to show the ability to speak the language “properly”.

Furthermore, register inflexibility was observed not only based on the general frequency distribution of post-predicate structures but also in the analysis of the information status of referents. The findings revealed that both heritage and majority Turkish speakers in Germany and Turkey place new referents in the post-predicate position, with heritage speakers doing so significantly more often. However, while majority speakers typically do this only in informal situations, heritage speakers employ this structure across all communicative contexts. As a result, the more frequent placement of new referents in the post-predicate position by heritage speakers can be attributed to register inflexibility.

As for Kurmanji, investigating language variation across communicative situations presents an interesting and innovative perspective. Unlike Turkish, there is no opposition between heritage and majority Kurmanji. Moreover, the language is used predominantly in informal situations, at least in Turkey. For instance, formal procedures like speaking to the police or writing a witness report are conducted in Turkish, not Kurmanji. Hence, Kurmanji-Turkish bilingual speakers are basically never exposed to Kurmanji in formal situations. However, they are likely aware of register variation in languages, owing to the knowledge they receive in Turkish. Thus, it is interesting to see whether speakers implement their register knowledge also in their heritage language, thereby creating novel strategies in the formal situations of Kurmanji.

Furthermore, as noted several times, no study has specifically investigated the use of the post-predicate position in Kurmanji across different communicative situations, at least to the best of our knowledge. Although research on Persian has demonstrated that a lower degree of formality corresponds with an increase in post-predicate elements (\citealt{Frommer1981}, \citealt{Haig_post-predicate_EtAl2024-1}).

The analysis of the Kurmanji data revealed that the frequency distribution of post-predicate structures showed no effect of communicative situations. This finding suggests that the phenomenon under investigation is not sensitive to formality or mode in Kurmanji, rather, it is linked to predicate semantics. However, in structures where previous research (\citealt{Haig2014, Gundogdu2019}) indicated potential variation—specifically regarding the position of goals marked with an adposition—the analysis found that formality does indeed influence the placement of such structures. Goals flagged with an adposition are favored in the post-predicate position in the informal situations. This preference for the post-predicate position mirrors patterns observed in Turkish. This suggests that when variation is possible in the heritage language, speakers may be inclined to adopt strategies known from the majority language.

Finally, the analysis of the Kurmanji data revealed that, in addition to goal arguments, non-goal constituents are also placed in the post-predicate position. An examination of the frequency of these constituents across communicative situations indicated that the spoken mode particularly facilitates their use in Kurmanji spoken in Turkey. A similar pattern has been observed in Turkish, where the spoken mode positively impacts the use of the post-predicate position. As previously mentioned, the spontaneous nature of spoken language allows for greater word order variation, whereas the written mode offers opportunities for corrections and adjustments. However, interpreting the results for the distribution of non-goal constituents across communicative situations should be approached with caution, as the total count of these structures is relatively low.

\section{Limitations of the study}\label{sec:key:7.4}

While the study provides an in-depth exploration of the topic, several limitations should be acknowledged. First of all, as emphasized several times, the data from the trilingual group in Germany comes from nine speakers. For this reason, the conclusions about this group are mainly drawn from qualitative patterns rather than statistical analyses. For future research, more data need to be collected from the trilingual group. Besides, another group could be added to the study, that is Kurmanji-German bilinguals, who do not have Turkish in their repertoire. Adding this group could shed light on the contact-induced language change in Kurmanji in Germany.

Another limitation, and at the same time a direction for future research, concerns the range of communicative situations investigated. First, the categorization of police reports as formal and messages to a friend as informal could be regarded as somewhat simplistic. In order to determine whether the findings can be generalized to formal and informal communicative situations more broadly, future research should replicate the study using a wider range of communicative contexts. Moreover, \citet{Maas2010} differentiates between informal registers in private and public contexts. An illustration of a private informal register is communication within family and friends, whereas a public informal register encompasses everyday interactions at workplaces, markets, streets, and similar settings. Thus, it would be interesting to investigate the use of the post-predicate position also in public informal registers. Besides, the study showed that the use of post-predicate structures in Turkish is also high in the formal communicative situation but only in the spoken mode. However, it is interesting to investigate whether such a finding is relevant for all formal spoken communicative situations or only for spontaneous formal speech. Hence, the investigation of the formal spoken communicative situations prepared in advance, such as a talk at a conference, a lecture, can be made.

Furthermore, for the analysis of the information status of post-predicate constituents, it is important to remember that the data from the four communicative situations were analyzed. Hence, the concept of ‘given’ and ‘new’ within one communicative situation might overlap in scenarios where a speaker has already narrated the story once, twice, or even more times. Thus, the information status of the constituents should be investigated with data that was not narrated several times. Another alternative is to use a different stimulus for the elicitation of different communicative situations.

Finally, while the analysis of post-predicate elements provides a detailed account of their weight, constituent types, semantic roles, the study does not examine the properties of constituents occurring in the pre-predicate position. Consequently, it is not possible to assess which types of constituents appear pre-predicatively or how their distribution relates to the patterns observed post-predicatively. Future research should therefore include a systematic analysis of constituents in the pre-predicate position in order to provide a more complete picture of constituent distribution.

